% =============================================================================
% Capítulo 6 – Plano de Trabalho e Cronograma
% =============================================================================
\chapter{Plano de Trabalho e Cronograma}
\label{cap_plano_trabalho}

O desenvolvimento do SIGEA-GTT seguiu a metodologia ágil adaptada
\textit{Scrum com Sprints} de duas semanas, complementada pelos princípios
do PDCA (\textit{Plan-Do-Check-Act}) para garantia da qualidade em cada
entrega. O projeto foi dividido em quatro Fases principais, cujo
planejamento temporal é apresentado no Diagrama de Gantt das
\autoref{fig_gantt1} e \autoref{fig_gantt2}.

\section{Descrição das Fases e Marcos}
\label{sec_fases_marcos}

\begin{table}[htb]
\centering
\small
\caption{Fases, marcos e entregas do projeto SIGEA-GTT.}
\label{tab_fases}
\begin{tabularx}{\textwidth}{c l X c}
\toprule
\textbf{Fase} & \textbf{Denominação} & \textbf{Entregas Principais} & \textbf{Período} \\
\midrule
1 & Fundação e Autenticação &
    Modelagem do banco (PostgreSQL 16 + Flyway), Spring Security 6 com JWT,
    CRUD de Usuários e Módulos/Gatilhos GTT, catálogo de Gravidades NCC MERP,
    tela de login e redefinição de senha obrigatória no Angular 18.
    & Mar – Abr/2025 \\
\midrule
2 & Gestão do Catálogo Clínico &
    Catálogo educacional de Módulos e Gatilhos GTT com \textit{accordions}
    e pistas clínicas; CRUD completo de Gravidades NCC MERP com guia
    interativo e distinção visual de dano real (E–I);
    testes unitários e de integração 100\% (JUnit 5 + Jacoco).
    & Mai – Jun/2025 \\
\midrule
3 & Módulo Educacional &
    CRUD de Turmas Acadêmicas, vinculação Professor/Alunos, Atividades com
    Prontuários Simulados (JSONB), cronômetro de 20 min, seis Ferramentas da
    Qualidade interativas (RF18–RF22), submissão, correção e Dashboard analítico
    IHI ($TI_{1000}$, $TI_{100}$, $P_{EA}$).
    & Jul – Set/2025 \\
\midrule
4 & Documentação Acadêmica &
    Relatório técnico ABNT (TCC DCOMP-UFS), seções de Engenharia de Requisitos,
    Design Thinking, diagramas TikZ (DER, Casos de Uso, Classes 4+1),
    Diagrama de Gantt, Apêndices e compilação final LaTeX.
    & Out/2025 – Jan/2026 \\
\bottomrule
\end{tabularx}
\fonte{Elaborado pelo autor (2026).}
\end{table}

\section{Diagrama de Gantt – Fase 1 e 2}
\label{sec_gantt_12}

\begin{figure}[htb]
\centering
\caption{Diagrama de Gantt – Fases 1 e 2 (março a junho de 2025).}
\label{fig_gantt1}

\resizebox{\textwidth}{!}{%
\begin{tikzpicture}
\begin{ganttchart}[
    hgrid,
    vgrid={*{3}{draw=none}, dotted},
    x unit=0.55cm,
    y unit title=0.7cm,
    y unit chart=0.65cm,
    title height=1,
    title label font=\scriptsize\bfseries,
    bar label font=\scriptsize,
    group label font=\scriptsize\bfseries,
    milestone label font=\scriptsize\itshape,
    bar/.style={fill=blue!50!black, rounded corners=2pt},
    bar incomplete/.style={fill=blue!20, rounded corners=2pt},
    group/.style={fill=blue!70, rounded corners=2pt},
    milestone/.style={fill=red!70, shape=diamond, inner sep=2pt},
    link/.style={->, thick, red!70!black},
    today=8, today rule/.style={draw=red, dashed, thick},
    today label={\scriptsize Hoje},
    calendar week text = \small{\startday},
]{1}{32}
  \gantttitle{Março 2025}{8}
  \gantttitle{Abril 2025}{8}
  \gantttitle{Maio 2025}{8}
  \gantttitle{Junho 2025}{8} \\

  % Fase 1
  \ganttgroup{FASE 1 – Fundação e Autenticação}{1}{16} \\
  \ganttbar{Modelagem BD + Flyway V1}{1}{4} \\
  \ganttbar{Spring Security 6 + JWT}{3}{8} \\
  \ganttbar{CRUD Usuários (Admin)}{5}{10} \\
  \ganttbar{CRUD Módulos GTT}{7}{12} \\
  \ganttbar{CRUD Gatilhos GTT}{9}{14} \\
  \ganttbar{Angular 18 – Login + 1.º Senha}{11}{16} \\
  \ganttmilestone{Marco 1 – Fase 1 Concluída}{16} \\[grid]

  % Fase 2
  \ganttgroup{FASE 2 – Gestão do Catálogo Clínico}{17}{32} \\
  \ganttbar{Catálogo GTT (Accordions)}{17}{22} \\
  \ganttbar{CRUD Gravidades NCC MERP}{19}{24} \\
  \ganttbar{Guia Clínico MERP (UI)}{22}{27} \\
  \ganttbar{Testes JUnit 5 + Jacoco}{25}{30} \\
  \ganttbar{Documentação Javadoc / OpenAPI}{27}{32} \\
  \ganttmilestone{Marco 2 – Fase 2 Concluída}{32} \\

  % Links
  \ganttlink{elem1}{elem2}
  \ganttlink{elem2}{elem3}
  \ganttlink{elem3}{elem4}
  \ganttlink{elem4}{elem5}
  \ganttlink{elem5}{elem6}
\end{ganttchart}
\end{tikzpicture}
}%
\fonte{Elaborado pelo autor com pgfgantt (2026).}
\end{figure}

\section{Diagrama de Gantt – Fase 3 e 4}
\label{sec_gantt_34}

\begin{figure}[htb]
\centering
\caption{Diagrama de Gantt – Fases 3 e 4 (julho de 2025 a janeiro de 2026).}
\label{fig_gantt2}

\resizebox{\textwidth}{!}{%
\begin{tikzpicture}
\begin{ganttchart}[
    hgrid,
    vgrid={*{3}{draw=none}, dotted},
    x unit=0.42cm,
    y unit title=0.7cm,
    y unit chart=0.65cm,
    title height=1,
    title label font=\scriptsize\bfseries,
    bar label font=\scriptsize,
    group label font=\scriptsize\bfseries,
    milestone label font=\scriptsize\itshape,
    bar/.style={fill=blue!50!black, rounded corners=2pt},
    bar incomplete/.style={fill=blue!20, rounded corners=2pt},
    group/.style={fill=green!50!black, rounded corners=2pt},
    milestone/.style={fill=red!70, shape=diamond, inner sep=2pt},
    link/.style={->, thick, red!70!black},
]{1}{52}
  \gantttitle{Jul/2025}{8}
  \gantttitle{Ago/2025}{8}
  \gantttitle{Set/2025}{8}
  \gantttitle{Out/2025}{8}
  \gantttitle{Nov/2025}{8}
  \gantttitle{Dez/2025}{8}
  \gantttitle{Jan/2026}{4} \\

  % Fase 3
  \ganttgroup{FASE 3 – Módulo Educacional}{1}{28} \\
  \ganttbar{CRUD Turmas Acadêmicas}{1}{6} \\
  \ganttbar{Atividades + Prontuário JSONB}{5}{12} \\
  \ganttbar{Cronômetro 20 min (Frontend)}{9}{14} \\
  \ganttbar{6 Ferramentas da Qualidade}{11}{20} \\
  \ganttbar{Submissão + Correção Docente}{17}{24} \\
  \ganttbar{Dashboard Analítico IHI}{21}{28} \\
  \ganttmilestone{Marco 3 – Fase 3 Concluída}{28} \\[grid]

  % Fase 4
  \ganttgroup{FASE 4 – Documentação Acadêmica}{28}{52} \\
  \ganttbar{Capítulos 1–2 (Intro/Metodologia)}{28}{33} \\
  \ganttbar{Capítulo 3 (Eng. Requisitos)}{32}{36} \\
  \ganttbar{Capítulo 4 (Design UX/Personas)}{35}{39} \\
  \ganttbar{Capítulo 5 (Diagramas TikZ)}{38}{43} \\
  \ganttbar{Capítulo 6 (Gantt)}{42}{45} \\
  \ganttbar{Capítulo 7 (Conclusão)}{44}{47} \\
  \ganttbar{Apêndices + Anexos}{46}{49} \\
  \ganttbar{Revisão ABNT + Compilação Final}{48}{52} \\
  \ganttmilestone{Marco 4 – TCC Concluído}{52} \\

  % Links Fase 3
  \ganttlink{elem1}{elem2}
  \ganttlink{elem2}{elem3}
  \ganttlink{elem3}{elem4}
  \ganttlink{elem4}{elem5}
  \ganttlink{elem5}{elem6}
\end{ganttchart}
\end{tikzpicture}
}%
\fonte{Elaborado pelo autor com pgfgantt (2026).}
\end{figure}

\section{Metodologia de Gestão e Controle de Qualidade}
\label{sec_gestao_qualidade}

A gestão das atividades de desenvolvimento foi realizada por meio de um
quadro Kanban digital (\textit{board} GitHub Projects), com as colunas:
\textbf{Backlog}, \textbf{Em Desenvolvimento}, \textbf{Em Revisão} e
\textbf{Concluído}. A política de finalização (\textit{Definition of Done})
exigiu, em cada tarefa:

\begin{enumerate}
    \item Implementação funcional completa e revisada pelo orientador;
    \item Cobertura de testes automatizados de 100\% de ramos e linhas
    (Jacoco no backend, Jest no frontend);
    \item Geração sem erros da documentação automática
    (Javadoc + SpringDoc OpenAPI 3 + Compodoc);
    \item Validação empírica da interface em todas as resoluções homologadas
    (HD, HD+, WUXGA e Ultrawide 21:9);
    \item Aprovação formal do orientador e da coorientadora antes da
    transição para a próxima fase.
\end{enumerate}
